\part{Introduction to LLM}

\section{What is a Language Model (LM)?}

A Language Model is a probabilistic model of a natural language. It produce a probability for any sequence of words.

The probability over all possible sentences in a language sums to 1.
There are region of high density in the space of all possible sentences.

\subsection{Formal definition for language models}

\begin{equation*}
  P(w1, w2, ..., wn) = \prod_{t=1}^{T} P(w_t | w_1, w_2, ..., w_{t-1})
\end{equation*}


\section{Simple LMs - n-gram models}

We need to estimate the probability of a word given the previous words.
A simplifying assumption, Markov assumption, is that:

\begin{equation*}
  P(w_t | w_1, w_2, ..., w_{t-1}) \approx P(w_t | w_{t-n+1}, ..., w_{t-1})
\end{equation*}

The next word $w_t$ depends only on a small windows of $n-1$ preceding words called the \textbf{context}.


For $n=2$, the next word only depends on the previous word, called a \textbf{bigram} model.
\begin{equation*}
  P(w_t | w_{t-1}, w_{t-2}, ..., w_1) \approx P(w_t | w_{t-1})
\end{equation*}


\section{Estimation \& generation}

We can use a language model to estimate probabilities for new
sentences.
Generate a new sentence based on known probabilities.

\begin{itemize}
  \item Starting from a given word, we can compute the probability of having any
        other word following it
  \item Then, we choose one word based on this probability distribution
  \item We add the new word to the already available sequence, and generate the
        next one
  \item We call this, autoregressive generation
\end{itemize}


\paragraph{Example}

$P(the, cat, chased, the, cheese) = P(the) \times P(cat | the) \times P(chased | the, cat) \times P(the | the, cat, chased) \times P(cheese | the, cat, chased, the)$

We can use the Markov assumption to simplify the computation of the probabilities.

$P(the, cat, chased, the, cheese) \approx P(the) \times P(cat | the) \times P(chased | cat) \times P(the | chased) \times P(cheese | the)$

Let's assume $P(the) = 1$

$P(the, cat, chased, the, cheese) \approx 1 \times \frac{1}{4} \times 1 \times 1 \times \frac{1}{4} = \frac{1}{16}$


\section{Generating new sentences}



\section{Limitation of n-grams}

\begin{enumerate}
  \item \textbf{Data sparsity}: As n increases, the number of possible n-grams grows exponentially. Bigram, V words, the table size is $V^2$. n-grams, V words, the table size is $V^n$. A small vocabulary or a small context size make the LM quite useless
  \item \textbf{Context limitation}: Only local context (previous n-1 words) are used to choose. Longer sentences may require remembering what happened ”early on”
  \item \textbf{Lack of semantics}:Similar words are treated in the same way as completely different ones
\end{enumerate}

The languafe models are not a new concept. In the '90s, researchers tried to build LMs based on existing
techniques and computing capabilities. With varying results... Mostly bad.


\section{Bigger and better Language Models}



We can see $P(w_t | w_{t-1}, w_{t-2}, ..., w_1)$ as a classification function.

"Given an \textbf{input}, the previous words in the sentence, predict the \textbf{next word}, class."



\chapter{Introduction to deep learning}

\section{The perceptron}

The perceptron is the simplest unit of neural networks. It takes an input and \textbf{multiplies} it by
a weight. Then, it sums all the weighted inputs and applies a non-linear function to produce an output.

\begin{equation*}
  y = f(w_0 + w_1x_1 + w_2x_2 + ... + w_nx_n)
\end{equation*}

\begin{figure}[h!]
  \centering
  \includegraphics[width=0.3\linewidth]{content/API_Programming/images/2026-09-21-09-33-07.png}
\end{figure}

Where:

\begin{itemize}
  \item $y$ is the output of the perceptron
  \item $f$ represents a non-linear “activation” function
  \item $w_i$ and $w_0$ are the weights and bias witch are learned during training
  \item $x = (x_1, x_2, ..., x_n)$ are the inputs to the perceptron
\end{itemize}


\section{The perceptron, in 2D}

\begin{figure}[h!]
  \centering
  \includegraphics[width=0.3\linewidth]{content/API_Programming/images/2026-09-21-09-36-40.png}
\end{figure}

\section{Activation functions}

Activation functions are used for two main reasons:

\begin{itemize}
  \item Enforce properties on perceptron's output. E.g. sigmoid -> binds output to [0, 1] range
  \item Introduce non-linearities in the model + some others (faster convergence, sparsity, ...)
\end{itemize}


Common activation functions are: ReLU, Sigmoid, Leaky ReLU, Tanh, Softmax, Linear, GeLU

\begin{figure}[h!]
  \centering
  \includegraphics[width=0.3\linewidth]{content/API_Programming/images/2026-09-21-09-41-37.png}
\end{figure}

\section{1. Enforce properties on perceptron's output}

Binary classification problem: seprate positive ($\star$) and negative ($\circ$) samples. For a point $x \in \R^2$,
the perceptron can predict $p(\star | x)$. For a binary case this implies that $p(\circ | x) = 1 - p(\star | x)$.

The Sigmoid maps any value in $\R$ to the range [0, 1]. The output is squashed between 0 and 1.

\begin{equation*}
  \sigma(x) = \frac{1}{1 + e^{-x}}
\end{equation*}

\subsection{Adding some perceptrons}

\begin{figure}[h!]
  \centering
  \includegraphics[width=0.3\linewidth]{content/API_Programming/images/2026-09-21-09-50-18.png}
\end{figure}


\subsection{and adding other layers!}

\begin{figure}[h!]
  \centering
  \includegraphics[width=0.3\linewidth]{content/API_Programming/images/2026-09-21-09-51-10.png}
\end{figure}

\section{2. Introduce non-linearities in the model}

This is the second reason why we need activation functions.

If $f(x) = x$ (i.e., no non-linearity is added) we get $z = s^T W^T x$


\section{Multi-layer perceptron models}



\section{Activation functions for classification models}

As argued, activation functions can be used to enforce properties on
the model's output

In classification problems, the output before the final activation is
treated as unnormalized probabilities \textbf{logits unit}.

We still need a step to convert logits into valid probabilities

\begin{figure}[h!]
  \centering
  \includegraphics[width=0.3\linewidth]{content/API_Programming/images/2026-09-21-10-18-55.png}
\end{figure}

\section{Binary classification}



\section{Multi-class classification}


The \textbf{softmax} function produce a probability of class (e.g. class 1 (cat), class 2 (dog)), it gives us
only values between 0 and 1

\begin{equation*}
  P(c_i|x) = \frac{e^{z_i}}{\sum_{j} e^{z_j}}
\end{equation*}

\begin{figure}[h!]
  \centering
  \includegraphics[width=0.3\linewidth]{content/API_Programming/images/2026-09-21-10-28-21.png}
\end{figure}

\section{Activation functions for regression models}



\section{Defining weights (parameters)}

So far, we assumed all weights and biases (let's call them $\theta$) to be
known. \textbf{But, we still need to figure out how we find them!}

We pick a function, onjective, or loss, $\L(\theta)$, that we want to minimize.

It looks at the model, the data are fixed, and give us a number, high number means bad model,
low number means good model.

\paragraph{Example:} $$\L(\theta) = MSE(\theta) = \frac{1}{n}\sum (y_i - \theta^T x_i)^2$$

Then, we pick $\theta$ that minimizes it

\section{Linear regression}

For simple models, we can find the optimal weights in closed form

\begin{equation*}
  \frac{\partial \L(\theta)}{\partial \theta} = \frac{\partial MSE(\Theta)}{\partial \theta} = 0
\end{equation*}

Quadratic in $\theta$, can be solved easily!

Or, we can evaluate the loss function for a bunch of $\theta$'s, and find the
“best” one

\begin{figure}[h!]
  \centering
  \includegraphics[width=0.3\linewidth]{content/API_Programming/images/2026-09-21-10-40-02.png}
\end{figure}

\section{More complex losses/model}

In more complex loss functions/models, we may not be able to solve the problem in a closed form.
But we can evaluate $\L(\theta)$  for various values of $\theta$.


We can iteratively update $\theta$ to reach a local minimum. We can start from a random value $\theta$, then:

\begin{itemize}
  \item We start from a random value $\theta$, then
  \item we “move around” according to “some policy”
\end{itemize}

\textbf{Move around }= update $\theta$ incrementally, based on its current value. The new value of $\theta$
at any step depends on the previous step's value: $\theta_{t+1} := \theta_t + \textit{update}$


\textbf{Some policy} = we take a small step in the direction where the
function decreases locally. i.e. in the \textbf{opposite} direction of the \textbf{gradient}.
$\theta_{t+1} := \theta_t - \alpha \nabla_\theta \L(\theta_t)$.

For 1-dimensional $\theta$, we have $\theta_{t+1} := \theta_t - \alpha \frac{\partial \L(\theta)}{\partial \theta}$
$\alpha$: learning rate, controls the “size” of the step ($10^{-3}$, $10^{-4}$, ...)

If the derivative is negative, i move to the right, when the derivative is positive i move to the left and
i get the bottom of the function.

Gradient Descent!

\subsection{Limitation of gradient descent}

GD is sensitive to weight initialization. Different initializations can lead to different solutions!.
GD can get stuck in local minima


Various solutions to help prevent local minima:
\begin{itemize}
  \item Adding momentum
  \item Adaptive learning rates, how to regulate the step
  \item Learning rate schedules
\end{itemize}

\section{Backpropagation}

How to compute the gradient of a loss function with respect to the weights of a neural network?

So far, we assumed we were able to compute $\nabla_\theta \L(\theta)$.

We need some automated way to compute gradients.

Now we introduce the \textbf{backpropagation}, to compute the gradient of the loss w.r.t.
any weight! Backpropagation is just a fancy word for “using the chain rule”

\subsection{Using the chain rule}

We use the chain rule from calculus

\begin{equation*}
  \frac{\partial f}{\partial x} = \frac{\partial f}{\partial g} \cdot \frac{\partial g}{\partial x}
\end{equation*}

Other notation caould be $f(g(x))' = f'(g(x)) \cdot g'(x)$

And apply it from the end of the computational graph, backwards. Hence the name, backpropagation.


\subsection{Computational graph}

A computational graph is a directed graph
\begin{itemize}
  \item Each node corresponds to an operation
  \item Each edge represents the flow of data between nodes
\end{itemize}

For instance, we may want to compute $y = wx + b$

\begin{figure}[h!]
  \centering
  \includegraphics[width=0.3\linewidth]{content/API_Programming/images/2026-09-21-11-02-42.png}
\end{figure}

We start from three variables, $w$, $x$ and $y$. The computational graph performs one
operation at a time.


\begin{itemize}
  \item First, compute the intermediate variable $a = wx$
  \item Then, compute the output variable $z = a + y = wx + y $
\end{itemize}

\paragraph{Backpropagation example}

Let's say:

\begin{itemize}
  \item Out dataset has one point ($x, y$)
  \item Our (weird) model has two parameters, $\theta_1$ and $\theta_2$, and predicts $\theta_1 \theta_2 x$
  \item Our loss function is $L = (\theta_1 \theta_2 x - y)^2$
\end{itemize}

We build a computational graph with all operations and intermediate
variables.

\begin{itemize}
  \item $ a = \theta_1 \theta_2$
  \item $ b = ax = \theta_1 \theta_2 x$
  \item $c = b - y =  ax - y = \theta_1 \theta_2 x - y$
  \item $L = c^2 = (b - y)^2 = (ax - y)^2 =  (\theta_1 \theta_2 x - y)^2$
\end{itemize}

\section{Loss functions}

What loss functin should I use?

\begin{itemize}
  \item Regression: Mean Squared Error (MSE), Mean Absolute Error (MAE)
  \item Binary: Binary Cross Entropy (BCE), the ground truth is $y = {0, 1}$ and the predicted value is $\hat{y} = model(x) \in [0, 1]$
\end{itemize}

\begin{equation*}
  \L = -y\log(\hat{y}) - (1-y)\log(1-\hat{y})
\end{equation*}

y, the ground truth, acts as a "selector" of the loss term to be applied.

\begin{itemize}
  \item $y = 1 \to \L = -\log(\hat{y})$
  \item $y = 0 \to \L = -\log(1-\hat{y})$
\end{itemize}



\begin{figure}[h!]
  \centering
  \includegraphics[width=0.3\linewidth]{content/API_Programming/images/2026-09-21-11-21-18.png}
\end{figure}


\section{Loss functions}

Multi-class classification

\begin{itemize}

  \item  Cross-Entropy
  \item  Generalization to multiple classes of BCE
  \item  $yF = 1$ when ground truth is the ith class, 0 otherwise
  \item  $yF$ plays the same “selector” mechanism as in BCE
\end{itemize}


\begin{equation*}
  \L = -\sum_{i} y_i \log(\hat{y}_i)
\end{equation*}

